\chapter{Topology}

Topology is the study of the qualitative structure of spaces. In Euclidean geometry, we measure length, angle, and area. In metric spaces, we measure distance. In topology, we keep only the information needed to speak about nearness, convergence, continuity, connectedness, and compactness.

This may seem like a loss of information, but it is in fact a powerful abstraction. Many theorems in analysis do not depend on the exact numerical value of a distance. The Intermediate Value Theorem depends on connectedness. The Extreme Value Theorem depends on compactness. Continuity can be described without mentioning a formula for distance. Topology isolates these structural ideas.

The basic object is a set together with a chosen collection of subsets called open sets. Once the open sets are specified, the language of limits and continuity can be rebuilt.

\section{From Metric Spaces to Topological Spaces}

\subsection{Open Sets in Metric Spaces}

Let us recall the metric-space intuition. If $(X,d)$ is a metric space, then an open ball centered at $x\in X$ with radius $r>0$ is
\[
    B_r(x)=\{y\in X:d(x,y)<r\}.
\]
A subset $U\subseteq X$ is open if for every $x\in U$, there exists $r>0$ such that $B_r(x)\subseteq U$.

\begin{example}
In $\mathbb{R}$ with the usual metric $d(x,y)=|x-y|$, the interval $(0,1)$ is open. For every $x\in(0,1)$, we may choose
\[
    r<\min\{x,1-x\},
\]
so that $(x-r,x+r)\subseteq(0,1)$.
\end{example}

The important observation is that once we know which subsets are open, we can define continuity without explicitly using the metric.

\begin{proposition}[Topological Form of Continuity in Metric Spaces]
Let $(X,d_X)$ and $(Y,d_Y)$ be metric spaces. A function $f:X\to Y$ is continuous if and only if for every open set $V\subseteq Y$, the preimage $f^{-1}(V)$ is open in $X$.
\end{proposition}

\begin{proof}
Assume first that $f$ is $\varepsilon$-$\delta$ continuous. Let $V\subseteq Y$ be open and let $x\in f^{-1}(V)$. Then $f(x)\in V$, so there exists $\varepsilon>0$ such that $B_\varepsilon(f(x))\subseteq V$. By continuity, there exists $\delta>0$ such that $d_X(x,x')<\delta$ implies $d_Y(f(x),f(x'))<\varepsilon$. Thus $B_\delta(x)\subseteq f^{-1}(V)$, so $f^{-1}(V)$ is open.

Conversely, assume preimages of open sets are open. For $x\in X$ and $\varepsilon>0$, the ball $B_\varepsilon(f(x))$ is open in $Y$, so $f^{-1}(B_\varepsilon(f(x)))$ is open in $X$ and contains $x$. Therefore it contains some ball $B_\delta(x)$, giving the usual $\varepsilon$-$\delta$ condition.
\end{proof}

This proposition motivates the axiomatic definition of topology.

\section{Topological Spaces}

\subsection{Definition of a Topology}

\begin{definition}[Topology]
Let $X$ be a set. A \textbf{topology} on $X$ is a collection $\mathcal{T}\subseteq\mathcal{P}(X)$ satisfying:
\begin{enumerate}
    \item $\emptyset\in\mathcal{T}$ and $X\in\mathcal{T}$;
    \item arbitrary unions of elements of $\mathcal{T}$ belong to $\mathcal{T}$;
    \item finite intersections of elements of $\mathcal{T}$ belong to $\mathcal{T}$.
\end{enumerate}
The pair $(X,\mathcal{T})$ is called a \textbf{topological space}. Elements of $\mathcal{T}$ are called \textbf{open sets}.
\end{definition}

The asymmetry between arbitrary unions and finite intersections is essential. In $\mathbb{R}$, arbitrary unions of open intervals are open, but arbitrary intersections of open sets need not be open.

\begin{example}
For each $n\in\mathbb{N}$, the interval $(-1/n,1/n)$ is open in $\mathbb{R}$. However,
\[
    \bigcap_{n=1}^{\infty}(-1/n,1/n)=\{0\},
\]
and $\{0\}$ is not open in the usual topology of $\mathbb{R}$.
\end{example}

\subsection{Fundamental Examples}

\begin{example}[Discrete Topology]
For any set $X$, the collection $\mathcal{T}=\mathcal{P}(X)$ is a topology. Every subset is open. This is called the \textbf{discrete topology}.
\end{example}

\begin{example}[Indiscrete Topology]
For any set $X$, the collection $\mathcal{T}=\{\emptyset,X\}$ is a topology. This is called the \textbf{indiscrete topology} or \textbf{trivial topology}.
\end{example}

\begin{example}[Usual Topology on $\mathbb{R}^n$]
The topology generated by Euclidean open balls in $\mathbb{R}^n$ is called the usual topology. Unless otherwise stated, $\mathbb{R}^n$ is equipped with this topology.
\end{example}

\begin{example}[Cofinite Topology]
Let $X$ be a set. Define
\[
    \mathcal{T}=\{\emptyset\}\cup\{U\subseteq X:X\setminus U\text{ is finite}\}.
\]
Then $\mathcal{T}$ is a topology on $X$, called the \textbf{cofinite topology}.
\end{example}

\begin{proof}
The empty set belongs to $\mathcal{T}$, and $X\setminus X=\emptyset$ is finite, so $X\in\mathcal{T}$. If $\{U_i\}_{i\in I}$ are non-empty open sets, then
\[
    X\setminus\bigcup_{i\in I}U_i=\bigcap_{i\in I}(X\setminus U_i),
\]
which is a subset of one finite set $X\setminus U_j$ if at least one $U_j$ is non-empty; hence it is finite. Thus arbitrary unions are open. If $U_1,\dots,U_n$ are open, then
\[
    X\setminus\bigcap_{k=1}^nU_k=\bigcup_{k=1}^n(X\setminus U_k),
\]
a finite union of finite sets, hence finite. Thus finite intersections are open.
\end{proof}

\begin{example}[Particular Point Topology]
Fix $p\in X$. Define
\[
    \mathcal{T}=\{\emptyset\}\cup\{U\subseteq X:p\in U\}.
\]
Then $\mathcal{T}$ is a topology. In this topology, every non-empty open set must contain the special point $p$.
\end{example}

\subsection{Comparing Topologies}

\begin{definition}[Finer and Coarser Topologies]
Let $\mathcal{T}_1$ and $\mathcal{T}_2$ be topologies on the same set $X$. We say $\mathcal{T}_1$ is \textbf{finer} than $\mathcal{T}_2$ if
\[
    \mathcal{T}_2\subseteq\mathcal{T}_1.
\]
Equivalently, $\mathcal{T}_2$ is \textbf{coarser} than $\mathcal{T}_1$.
\end{definition}

A finer topology has more open sets. More open sets make it easier for subsets to be open, but harder for functions into the space to be continuous.

\begin{example}
On any set $X$, the discrete topology is the finest topology, while the indiscrete topology is the coarsest topology.
\end{example}

\begin{remark}
Topology is not only about sets; it is about sets together with a chosen standard of local visibility. A subset is open if the topology allows us to see it as a neighborhood-like object.
\end{remark}

\section{Bases and Subbases}

It is often inconvenient to list every open set. Instead, we describe a smaller collection of basic open sets from which all open sets can be built.

\subsection{Basis for a Topology}

\begin{definition}[Basis]
Let $X$ be a set. A collection $\mathcal{B}\subseteq\mathcal{P}(X)$ is a \textbf{basis} for a topology on $X$ if:
\begin{enumerate}
    \item for every $x\in X$, there exists $B\in\mathcal{B}$ such that $x\in B$;
    \item if $x\in B_1\cap B_2$ with $B_1,B_2\in\mathcal{B}$, then there exists $B_3\in\mathcal{B}$ such that
    \[
        x\in B_3\subseteq B_1\cap B_2.
    \]
\end{enumerate}
The topology generated by $\mathcal{B}$ consists of all unions of elements of $\mathcal{B}$.
\end{definition}

\begin{example}
Open intervals $(a,b)$ form a basis for the usual topology on $\mathbb{R}$. Open balls form a basis for the usual topology on $\mathbb{R}^n$.
\end{example}

\begin{proposition}
If $\mathcal{B}$ is a basis on $X$, then the collection of all unions of basis elements is a topology on $X$.
\end{proposition}

\begin{proof}
The empty union gives $\emptyset$, and the first basis condition ensures that $X$ is the union of all basis elements. Arbitrary unions of unions of basis elements are still unions of basis elements. For finite intersections, it is enough to consider two open sets $U$ and $V$. If $x\in U\cap V$, then $x\in B_1\subseteq U$ and $x\in B_2\subseteq V$ for some basis elements. By the second basis condition, there exists $B_3$ such that $x\in B_3\subseteq B_1\cap B_2\subseteq U\cap V$. Thus $U\cap V$ is a union of basis elements.
\end{proof}

\subsection{Subbasis}

\begin{definition}[Subbasis]
A collection $\mathcal{S}\subseteq\mathcal{P}(X)$ is a \textbf{subbasis} for a topology on $X$ if the set of all finite intersections of elements of $\mathcal{S}$ forms a basis. The topology generated by $\mathcal{S}$ is the collection of arbitrary unions of such finite intersections.
\end{definition}

\begin{example}
In $\mathbb{R}$, the collection of all rays $(-\infty,a)$ and $(a,\infty)$ is a subbasis for the usual topology.
\end{example}

\subsection{Countability Axioms}

\begin{definition}[First and Second Countability]
A topological space $X$ is \textbf{first countable} if every point has a countable local basis. It is \textbf{second countable} if the topology has a countable basis.
\end{definition}

\begin{example}
The usual topology on $\mathbb{R}^n$ is second countable. A countable basis is given by all open balls whose centers have rational coordinates and whose radii are positive rational numbers.
\end{example}

\begin{proposition}
Every second countable space is first countable.
\end{proposition}

\begin{proof}
Let $\mathcal{B}$ be a countable basis for $X$. For a fixed point $x\in X$, the collection of all basis elements containing $x$ is countable. It is a local basis at $x$.
\end{proof}

\section{Closed Sets, Closure, Interior, and Boundary}

Open sets are primary in the axioms, but closed sets are equally important.

\begin{definition}[Closed Set]
Let $X$ be a topological space. A subset $C\subseteq X$ is \textbf{closed} if $X\setminus C$ is open.
\end{definition}

\begin{proposition}[Properties of Closed Sets]
In a topological space $X$:
\begin{enumerate}
    \item $\emptyset$ and $X$ are closed;
    \item arbitrary intersections of closed sets are closed;
    \item finite unions of closed sets are closed.
\end{enumerate}
\end{proposition}

\begin{proof}
These properties follow from the axioms for open sets and De Morgan's laws:
\[
    X\setminus\bigcap_i C_i=\bigcup_i(X\setminus C_i),
    \qquad
    X\setminus\bigcup_{k=1}^n C_k=\bigcap_{k=1}^n(X\setminus C_k).
\]
\end{proof}

\subsection{Closure and Interior}

\begin{definition}[Closure]
Let $A\subseteq X$. The \textbf{closure} of $A$, denoted $\overline{A}$, is the intersection of all closed sets containing $A$.
\end{definition}

\begin{definition}[Interior]
Let $A\subseteq X$. The \textbf{interior} of $A$, denoted $A^\circ$ or $\operatorname{Int}(A)$, is the union of all open sets contained in $A$.
\end{definition}

\begin{definition}[Boundary]
The \textbf{boundary} of $A$ is
\[
    \partial A=\overline{A}\setminus A^\circ.
\]
Equivalently,
\[
    \partial A=\overline{A}\cap\overline{X\setminus A}.
\]
\end{definition}

\newpage

\begin{example}
In $\mathbb{R}$ with the usual topology, for $A=(0,1)$,
\[
    \overline{A}=[0,1],
    \qquad
    A^\circ=(0,1),
    \qquad
    \partial A=\{0,1\}.
\]
For $A=\mathbb{Q}\subseteq\mathbb{R}$,
\[
    \overline{\mathbb{Q}}=\mathbb{R},
    \qquad
    \operatorname{Int}(\mathbb{Q})=\emptyset.
\]
\end{example}

\begin{proposition}[Closure via Open Neighborhoods]
Let $A\subseteq X$ and $x\in X$. Then $x\in\overline{A}$ if and only if every open set $U$ containing $x$ intersects $A$.
\end{proposition}

\begin{proof}
If some open neighborhood $U$ of $x$ is disjoint from $A$, then $X\setminus U$ is a closed set containing $A$ but not containing $x$, so $x\notin\overline{A}$. Conversely, if $x\notin\overline{A}$, then $X\setminus\overline{A}$ is an open neighborhood of $x$ disjoint from $A$.
\end{proof}

\subsection{Limit Points and Density}

\begin{definition}[Limit Point]
Let $A\subseteq X$. A point $x\in X$ is a \textbf{limit point} of $A$ if every open neighborhood of $x$ contains a point of $A$ different from $x$.
\end{definition}

\begin{definition}[Dense Subset]
A subset $A\subseteq X$ is \textbf{dense} in $X$ if $\overline{A}=X$.
\end{definition}

\begin{example}
Both $\mathbb{Q}$ and $\mathbb{R}\setminus\mathbb{Q}$ are dense in $\mathbb{R}$ with the usual topology. Every open interval contains both rational and irrational numbers.
\end{example}

\section{Constructing New Spaces}

Topology becomes powerful when we can construct new spaces from old spaces.

\subsection{Subspace Topology}

\begin{definition}[Subspace Topology]
Let $(X,\mathcal{T})$ be a topological space and let $Y\subseteq X$. The \textbf{subspace topology} on $Y$ is
\[
    \mathcal{T}_Y=\{Y\cap U:U\in\mathcal{T}\}.
\]
\end{definition}

\begin{example}
The interval $[0,1]$ as a subspace of $\mathbb{R}$ has open sets of the form $[0,1]\cap U$, where $U$ is open in $\mathbb{R}$. For example, $[0,1/2)$ is open in $[0,1]$, although it is not open in $\mathbb{R}$.
\end{example}

\subsection{Product Topology}

\begin{definition}[Product Topology]
Let $X$ and $Y$ be topological spaces. The \textbf{product topology} on $X\times Y$ is generated by the basis
\[
    \mathcal{B}=\{U\times V:U\subseteq X\text{ open},\ V\subseteq Y\text{ open}\}.
\]
\end{definition}

\begin{example}
The usual topology on $\mathbb{R}^2$ is the same as the product topology on $\mathbb{R}\times\mathbb{R}$, where each copy of $\mathbb{R}$ has the usual topology.
\end{example}

\begin{proposition}[Continuity into a Product]
A map $f:Z\to X\times Y$ is continuous if and only if its coordinate maps $\pi_X\circ f:Z\to X$ and $\pi_Y\circ f:Z\to Y$ are continuous.
\end{proposition}

\subsection{Quotient Topology}

\begin{definition}[Quotient Topology]
Let $X$ be a topological space, let $Y$ be a set, and let $q:X\to Y$ be a surjective map. The \textbf{quotient topology} on $Y$ is defined by declaring $U\subseteq Y$ open if and only if $q^{-1}(U)$ is open in $X$.
\end{definition}

\begin{example}[Circle as a Quotient]
The circle $S^1$ can be obtained by identifying the endpoints of the interval $[0,1]$. Formally, define an equivalence relation on $[0,1]$ by $0\sim 1$ and $x\sim x$ otherwise. The quotient space $[0,1]/\sim$ is homeomorphic to $S^1$.
\end{example}

\begin{remark}
Quotient topology is the rigorous language of gluing. Many geometric spaces are built by taking a simple space and identifying selected points or subspaces.
\end{remark}

\section{Separation Axioms}

Topological spaces can be very strange. Separation axioms measure how well a topology distinguishes points.

\begin{definition}[$T_0$, $T_1$, and Hausdorff Spaces]
Let $X$ be a topological space.
\begin{enumerate}
    \item $X$ is $T_0$ if for any distinct $x,y\in X$, there exists an open set containing one of them but not the other.
    \item $X$ is $T_1$ if for any distinct $x,y\in X$, there exists an open set containing $x$ but not $y$, and an open set containing $y$ but not $x$.
    \item $X$ is \textbf{Hausdorff} or $T_2$ if for any distinct $x,y\in X$, there exist disjoint open sets $U,V$ such that $x\in U$ and $y\in V$.
\end{enumerate}
\end{definition}

\begin{proposition}
Every Hausdorff space is $T_1$, and every $T_1$ space is $T_0$.
\end{proposition}

\begin{example}
Every metric space is Hausdorff. If $x\neq y$, let $r=d(x,y)/3$. Then $B_r(x)$ and $B_r(y)$ are disjoint open neighborhoods.
\end{example}

\begin{proposition}
A space $X$ is $T_1$ if and only if every singleton set $\{x\}$ is closed.
\end{proposition}

\begin{proof}
If $X$ is $T_1$, then for every $y\neq x$, there is an open set $U_y$ containing $y$ but not $x$. Thus
\[
    X\setminus\{x\}=\bigcup_{y\neq x}U_y
\]
is open, so $\{x\}$ is closed. Conversely, if every singleton is closed, then $X\setminus\{y\}$ is an open set containing $x$ but not $y$ whenever $x\neq y$.
\end{proof}

\begin{theorem}[Uniqueness of Limits in Hausdorff Spaces]
If $X$ is Hausdorff, then a sequence in $X$ has at most one limit.
\end{theorem}

\begin{proof}
Suppose $x_n\to x$ and $x_n\to y$ with $x\neq y$. Since $X$ is Hausdorff, choose disjoint open sets $U$ and $V$ with $x\in U$ and $y\in V$. Eventually $x_n\in U$, and eventually $x_n\in V$. Hence eventually $x_n\in U\cap V$, impossible because $U\cap V=\emptyset$.
\end{proof}

\section{Continuity and Homeomorphism}

\subsection{Continuous Functions}

\begin{definition}[Continuity]
Let $X$ and $Y$ be topological spaces. A function $f:X\to Y$ is \textbf{continuous} if for every open set $V\subseteq Y$, the preimage $f^{-1}(V)$ is open in $X$.
\end{definition}

\begin{proposition}[Equivalent Form Using Closed Sets]
A function $f:X\to Y$ is continuous if and only if for every closed set $C\subseteq Y$, the preimage $f^{-1}(C)$ is closed in $X$.
\end{proposition}

\begin{proof}
This follows from
\[
    f^{-1}(Y\setminus C)=X\setminus f^{-1}(C).
\]
\end{proof}

\begin{proposition}[Composition of Continuous Functions]
If $f:X\to Y$ and $g:Y\to Z$ are continuous, then $g\circ f:X\to Z$ is continuous.
\end{proposition}

\begin{proof}
For any open set $W\subseteq Z$,
\[
    (g\circ f)^{-1}(W)=f^{-1}(g^{-1}(W)).
\]
Since $g$ is continuous, $g^{-1}(W)$ is open in $Y$. Since $f$ is continuous, $f^{-1}(g^{-1}(W))$ is open in $X$.
\end{proof}

\subsection{Homeomorphisms}

\begin{definition}[Homeomorphism]
A function $f:X\to Y$ is a \textbf{homeomorphism} if:
\begin{enumerate}
    \item $f$ is bijective;
    \item $f$ is continuous;
    \item $f^{-1}$ is continuous.
\end{enumerate}
If such a map exists, $X$ and $Y$ are called \textbf{homeomorphic}.
\end{definition}

Homeomorphic spaces are considered the same from the topological viewpoint. They may look geometrically different, but their open-set structures are equivalent.

\begin{example}
The open intervals $(0,1)$ and $\mathbb{R}$ are homeomorphic. One possible homeomorphism is
\[
    f(x)=\tan\left(\pi x-\frac{\pi}{2}\right).
\]
\end{example}

\begin{remark}
A bijective continuous function need not be a homeomorphism. The inverse must also be continuous. However, a continuous bijection from a compact space to a Hausdorff space is always a homeomorphism.
\end{remark}

\section{Connectedness}

Connectedness formalizes the idea that a space consists of one piece.

\subsection{Definitions}

\begin{definition}[Separation]
A \textbf{separation} of a topological space $X$ is a pair of non-empty disjoint open sets $U,V\subseteq X$ such that
\[
    X=U\cup V.
\]
\end{definition}

\begin{definition}[Connected Space]
A topological space $X$ is \textbf{connected} if it has no separation.
\end{definition}

\begin{proposition}[Equivalent Characterization]
A space $X$ is connected if and only if the only subsets of $X$ that are both open and closed are $\emptyset$ and $X$.
\end{proposition}

\begin{proof}
If $A\subseteq X$ is both open and closed and is neither $\emptyset$ nor $X$, then $A$ and $X\setminus A$ form a separation. Conversely, if $X=U\cup V$ is a separation, then $U$ is open and also closed because $U=X\setminus V$.
\end{proof}

\begin{theorem}[Intervals Are Connected]
Every interval in $\mathbb{R}$ is connected.
\end{theorem}

\begin{proof}
We prove the idea for an interval $I$. Suppose $I=A\cup B$ is a separation, with $a\in A$ and $b\in B$, say $a<b$. Let
\[
    S=A\cap [a,b].
\]
Using the least upper bound property of $\mathbb{R}$, one shows that $c=\sup S$ must belong neither consistently to $A$ nor to $B$ without contradicting the openness of $A$ and $B$ in the subspace $I$. Hence no such separation exists.
\end{proof}

\subsection{Continuous Images of Connected Spaces}

\begin{theorem}
If $f:X\to Y$ is continuous and $X$ is connected, then $f(X)$ is connected.
\end{theorem}

\begin{proof}
Suppose $f(X)$ were disconnected. Then $f(X)=U\cup V$ for non-empty disjoint open sets in the subspace topology of $f(X)$. The preimages $f^{-1}(U)$ and $f^{-1}(V)$ are non-empty disjoint open subsets of $X$ whose union is $X$, contradicting connectedness.
\end{proof}

\begin{theorem}[Intermediate Value Theorem]
Let $f:[a,b]\to\mathbb{R}$ be continuous. If $f(a)<y<f(b)$ or $f(b)<y<f(a)$, then there exists $c\in(a,b)$ such that $f(c)=y$.
\end{theorem}

\begin{proof}
The interval $[a,b]$ is connected, and the continuous image $f([a,b])$ is connected in $\mathbb{R}$. Connected subsets of $\mathbb{R}$ are intervals. Hence $f([a,b])$ contains every value between $f(a)$ and $f(b)$.
\end{proof}

\subsection{Path Connectedness}

\begin{definition}[Path Connectedness]
A space $X$ is \textbf{path connected} if for any $x,y\in X$, there exists a continuous map $\gamma:[0,1]\to X$ such that $\gamma(0)=x$ and $\gamma(1)=y$.
\end{definition}

\begin{proposition}
Every path connected space is connected.
\end{proposition}

\begin{proof}
Suppose $X$ is path connected but disconnected, with separation $X=U\cup V$. Choose $x\in U$ and $y\in V$. Let $\gamma:[0,1]\to X$ be a path from $x$ to $y$. Then $\gamma([0,1])$ is connected as a continuous image of a connected interval, but it intersects both $U$ and $V$, giving a separation of $\gamma([0,1])$. This is impossible.
\end{proof}

\begin{remark}
The converse is false. There exist connected spaces that are not path connected, such as the topologist's sine curve. Thus path connectedness is stronger than connectedness.
\end{remark}

\section{Compactness}

Compactness is one of the central ideas of topology and analysis. In finite sets, many arguments work because a finite cover has finitely many pieces. Compactness generalizes this finite behavior to infinite spaces.

\subsection{Open Covers}

\begin{definition}[Open Cover]
Let $X$ be a topological space and $A\subseteq X$. An \textbf{open cover} of $A$ is a collection of open sets $\{U_i\}_{i\in I}$ such that
\[
    A\subseteq\bigcup_{i\in I}U_i.
\]
A \textbf{subcover} is a subcollection that still covers $A$.
\end{definition}

\begin{definition}[Compactness]
A topological space $X$ is \textbf{compact} if every open cover of $X$ has a finite subcover.
\end{definition}

\begin{example}
Every finite topological space is compact. Given any open cover, choose one open set containing each point. Since there are finitely many points, this gives a finite subcover.
\end{example}

\begin{example}
The open interval $(0,1)$ is not compact in the usual topology. The collection
\[
    \left\{\left(\frac{1}{n},1\right):n\geq 2\right\}
\]
is an open cover of $(0,1)$, but no finite subcollection covers points sufficiently close to $0$.
\end{example}

\subsection{Compactness in $\mathbb{R}^n$}

\begin{theorem}[Heine-Borel Theorem]
A subset $K\subseteq\mathbb{R}^n$ is compact if and only if it is closed and bounded.
\end{theorem}

\begin{remark}
The Heine-Borel theorem is special to finite-dimensional Euclidean spaces. In general metric spaces, closed and bounded does not necessarily imply compact.
\end{remark}

\begin{theorem}[Bolzano-Weierstrass Theorem]
Every bounded sequence in $\mathbb{R}^n$ has a convergent subsequence.
\end{theorem}

\begin{remark}
In $\mathbb{R}^n$, compactness, sequential compactness, and closed boundedness are equivalent. In arbitrary topological spaces, these equivalences may fail.
\end{remark}

\subsection{Continuous Images of Compact Spaces}

\begin{theorem}
If $f:X\to Y$ is continuous and $X$ is compact, then $f(X)$ is compact.
\end{theorem}

\begin{proof}
Let $\{V_i\}_{i\in I}$ be an open cover of $f(X)$. Then $\{f^{-1}(V_i)\}_{i\in I}$ is an open cover of $X$. Since $X$ is compact, finitely many preimages cover $X$, say $f^{-1}(V_{i_1}),\dots,f^{-1}(V_{i_n})$. Then $V_{i_1},\dots,V_{i_n}$ cover $f(X)$.
\end{proof}

\begin{theorem}[Extreme Value Theorem]
Let $X$ be compact and let $f:X\to\mathbb{R}$ be continuous. Then $f$ attains a maximum and a minimum on $X$.
\end{theorem}

\begin{proof}
The image $f(X)$ is compact in $\mathbb{R}$. By Heine-Borel, it is closed and bounded. Since it is bounded, it has a supremum and an infimum. Since it is closed, these values belong to $f(X)$. Therefore $f$ attains its maximum and minimum.
\end{proof}

\begin{theorem}[Compact to Hausdorff]
Let $X$ be compact, $Y$ Hausdorff, and $f:X\to Y$ a continuous bijection. Then $f$ is a homeomorphism.
\end{theorem}

\begin{proof}
It suffices to show that $f$ is a closed map. Let $C\subseteq X$ be closed. Since $X$ is compact, $C$ is compact. The continuous image $f(C)$ is compact in $Y$. In a Hausdorff space, compact subsets are closed. Hence $f(C)$ is closed, so $f^{-1}$ is continuous.
\end{proof}

\section{Metric Spaces Revisited}

Many topological spaces come from metrics, but not all topologies are metrizable.

\begin{definition}[Metrizable Space]
A topological space $(X,\mathcal{T})$ is \textbf{metrizable} if there exists a metric $d$ on $X$ such that $\mathcal{T}$ is exactly the topology generated by the open balls of $d$.
\end{definition}

\newpage

\begin{example}
Every discrete topological space is metrizable. Define
\[
    d(x,y)=
    \begin{cases}
    0,&x=y,\\
    1,&x\neq y.
    \end{cases}
\]
Then every singleton is open, so every subset is open.
\end{example}

\begin{remark}
Metrizability is a strong property. Metric spaces are always Hausdorff and first countable. Therefore any topology failing these properties cannot come from a metric.
\end{remark}

\section{Conclusion}

Topology begins by replacing distance with open sets. From this simple abstraction, we recover the language of continuity, convergence, compactness, and connectedness. The power of topology is that it reveals which theorems of analysis depend on metric quantities and which depend only on the qualitative structure of open sets.

The essential pattern is the same as in previous chapters: identify the primitive data, state axioms, define structure-preserving maps, and then study properties invariant under those maps. In topology, the primitive data are open sets, the structure-preserving maps are continuous functions, and the notion of sameness is homeomorphism.
